% Oscillations dans un circuit (R, L, C).
Une autre solution pour déterminer la valeur de la capacité du condensateur est
d'établir des oscillations électriques dans un circuit ($R, L, C$). Le
condensateur, préalablement chargé sous une tension $E=\qty{5.0}{\V}$, est
relié à une bobine d'inductance $L=\qty{1.0}{\H}$ et de résistance
$r=\qty{20}{\ohm}$, selon le schéma suivant~:

\begin{figure}[H]
    \centering
    \begin{circuitikz}
        \ctikzset{bipoles/cuteswitch/height=0.4}
        \newdimen\d \d=4cm
        \draw (0,0)
            to[C,l_=$C$,v^>=$u$] ++(0,\d)
            to[cute open switch,l=$\mathrm{K}$] ++(0.8\d,0)
            to[L,l={($L=\qty{1}{\H}$, $r=\qty{20}{\ohm}$)}] ++(0,-\d)
            to[short] (0,0);
   \end{circuitikz}
\end{figure}

L'acquisition de la tension aux bornes du condensateur permet d’obtenir la
courbe suivante~:

\begin{figure}[H]
    \centering
    \includegraphics[width=.6\textwidth]{polynesie_2010_ex2_8}
\end{figure}

\begin{enumerate}
    \item Qualifier le régime d'oscillations obtenu.
    \item Déterminer la valeur d'une grandeur temporelle liée aux
          oscillations. Donner son nom.
    \item Donner l'expression de la période propre $T_{0}$ de l'oscillateur en
          fonction de $L$ et de $C$.
    \item En assimilant la grandeur temporelle précédente à cette valeur, en
          déduire la capacité du condensateur.
\end{enumerate}
